% Kapitel 5 — Budget: 16 Seiten — Schreiben: 26.10.–27.11. (5.1–5.2 bis 13.11., 5.3–5.6 16.–27.11.)
% Vorgehen und Entwicklung; Ergebnisdarstellung in Kap. 7 (Beschluss 26.10.),
% volle Tabellen in den Anhang.

\chapter{Vormessungen}
\label{chap:vormessungen}

\section{Wie gemessen wird}% 3 S.
% Inhalt: Aufgaben (SWE-bench, Gold gleich Patch), Bedingungen, Messgrößen,
% Statistik (Vorzeichentest, Bootstrap), Vorregistrierung und Holdout, Nachweis
% (Messbäume, Pfad-Audit, eingefrorenes Binary, Provenance).
\label{sec:wie-gemessen}

\section{Entwicklung des Abrufs auf SWE-bench}% 3 S.
% Inhalt: Seed-Fix ausführlich, die übrigen Messungen je ein bis drei Sätze
% (Frage, Antwort, Änderung am Werkzeug).
\label{sec:entwicklung-swebench}

\section{Agentenverhalten (AP8)}% 3 S.
% Inhalt: dazu Agentenläufe auf teamable 05.10. und 09.10. als Entwicklungsmessung.
\label{sec:agentenverhalten}

\section{Ferne Aufgaben: Hard-Slice}% 3 S.
% Inhalt: Bestätigungsmessung, alter Default.
\label{sec:hard-slice}

\section{Seed-Änderung und Tests als Gold}% 2 S.
% Inhalt: Entwicklungswert, neuer Default, Stand der Studie.
\label{sec:seed-aenderung}

\section{Zusammenhänge und Überleitung zur Studie}% 2 S.
% Inhalt: Bilanz, was die Messungen nicht beantworten, Begründung des Wechsels
% der Hauptfrage.
\label{sec:ueberleitung-studie}
